% ===== §21 The Political Economy of the Field =====
\section{The Political Economy of the Field}
\label{sec:political-economy}

\noindent
This section is the one the whole paper has been pointing toward and deferring, the place where the account must answer for its politics or be convicted of the very thing it has repeatedly guarded against. The charge is the oldest one against any aesthetics of life: that it is a doctrine for those with the leisure to practise it, a refinement available to the unhurried and irrelevant or insulting to those whose days are consumed by labour that is not theirs to shape. An aesthetic education of the everyday is peculiarly exposed to this charge, because the everyday it addresses is, for many, precisely the site of exhaustion, the repeated labour performed for another's accumulation, and to bring such people a doctrine of building beautiful fields could be the cruelest possible irrelevance. The section meets the charge directly, and its answer is not that the charge is mistaken but that the account, properly understood, is the charge's refutation rather than its instance, because the cultivation it describes is the one form of aesthetic life that does not presuppose the leisure the charge names, and the people the charge invokes are exactly those the account most concerns.

\subsection{The charge at full strength}
\label{sub:pe-charge}

The charge must be stated at full strength, because a weak version would be easy to dismiss and the easy dismissal would be an evasion. Here it is. The building of a generative field, as this paper describes it, takes time, attention, and a measure of control over one's own days; it asks one to hold a half-second before a habitual act, to let contingency revise the morning, to return to the same daily thing and find it thickened, and all of this presupposes that one's days are one's own to shape, that the morning is not already wholly claimed by labour performed under another's clock, that the attention is not already exhausted by the demands of survival. For those whose every hour is sold, whose repetition is the alienated repetition whose yield is captured, whose exhaustion at the end of the day's labour leaves nothing for the patient cultivation of a field, the doctrine of aesthetic education is not false but inaccessible, a luxury of those who have the surplus of time and attention it silently requires. And to such people the offer of this doctrine is worse than useless, because it adds to their dispossession the suggestion that their failure to live beautifully is a failure of cultivation rather than a theft of the conditions cultivation requires.

This is the charge, and the paper does not soften it, because the conditions it names are real and the dispossession it describes is the actual situation of many. The political economy of the field is unequal: the time, attention, and self-direction that field-building draws on are distributed unequally, captured disproportionately from those whose labour is owned, and an account that pretended otherwise, that spoke of building beautiful fields as though everyone stood equally able to build them, would be ideology in the exact sense, a representation of a class-bound possibility as a universal one. The paper has insisted throughout on the orthogonality of refinement and justice, and here that orthogonality returns as a question about the account itself: a doctrine of refinement that ignored the justice of who can afford to be refined would be the fine and the unjust at the level of theory, an aestheticism that mistook its own class conditions for the human condition.

\subsection{The account refutes the charge}
\label{sub:pe-refutation}

The answer turns on a feature of the account built in from the start, deliberately and against the grain of every aesthetics of leisure, the feature that the cultivation it describes requires no surplus of time or means but only a different way of entering the hours one already has. The four practices, as the worked example set out, are not added exercises requiring spare hours; each is a re-meeting of an act the day already contains, the half-second held before a motion the day demanded anyway, the contingency received in an event that occurred regardless, the brief sharing of a moment that passed in any case, the return to a daily thing one was going to do again. The cultivation does not ask for hours the dispossessed lack; it asks for a changed manner of living the hours there are, and those hours, however owned, however exhausting, still contain the acts that the four practices re-meet. This was not an incidental convenience of the account but its deliberate construction, the reason the paper located aesthetic education in the everyday and in the most repetitive and unspacious everyday at that, because it is precisely the life with no spare hours that needed an aesthetic education making no demand for them.

\claim{The account refutes the charge of leisure-dependence rather than instancing it, because the cultivation it describes requires no surplus of time or means but only a changed manner of entering the hours one already has. The four practices re-meet acts the day already contains and add none; they ask not for spare hours but for a different way of living the hours there are. The aesthetic education of the everyday is therefore available precisely to those whose days hold no surplus, which is why it was located in the everyday and not in the exceptional, and why the people the charge invokes, the dispossessed whose repetition is exhausting, are the ones the account most concerns rather than the ones it excludes.}

This does not dissolve the inequality, and the section must be careful not to let the refutation overreach into a false consolation. The inequality is real: exhaustion is real, and a body and attention worn down by owned labour have less to bring even to a cultivation that asks for no spare hours, so that the dispossessed are not equally placed even with respect to a practice that requires no surplus. What the account establishes is the weaker and the true claim, that the cultivation does not presuppose leisure and so is not in principle barred to those without it, that the door is not closed by the lack of surplus hours, and this is enough to refute the charge that aesthetic education is by nature a doctrine of the leisured, while it is not enough to pretend that the worn-down come to the door as easily as the unhurried. The honest position holds both: the cultivation is available to the dispossessed, requiring no surplus they lack; and the dispossession still weighs, the exhaustion still costs, so that to make the cultivation truly available to all is not only to design it without a demand for leisure, which the paper has done, but to change the conditions that exhaust, which the paper cannot do and does not pretend to.

\subsection{The structural remainder}
\label{sub:pe-structural}

What the cultivation cannot reach is named, finally, as the structural remainder, the part of the problem that no aesthetic education can solve because it is not aesthetic. The capture of the everyday's yield, the ownership of the repetition by an order that takes its accumulation away, the exhaustion manufactured by labour performed under another's clock, these are structural conditions, and their remedy is structural, a change in who owns the yield of the cycle and who commands the hours of the day. The beauty of maintenance presupposed the return of the cycle's yield to the field; the political economy is the examination of where that presupposition fails, and where it fails the failure is not to be met with a finer cultivation but with a change in the structure that captures the yield. This is the same insistence the paper has made at every turn where the aesthetic met the structural, the refusal to aestheticize across a structural line, and it returns here as the account's final honesty about its own limits. Aesthetic education is necessary and it is available to the dispossessed, requiring no leisure they lack; and it is not sufficient, because the conditions that exhaust the dispossessed and capture their yield are structural conditions that only structural change can remedy. The paper offers the cultivation as what it is, a real good genuinely available in the unspacious everyday, and it declines to offer it as what it is not, a substitute for the structural change that the justice of the field, as distinct from its beauty, finally requires.

\begin{casebox}[Two yields, one captured and one not]
A worker comes home from a day whose every hour was sold, whose repetition built a wealth that flowed entirely to an owner, whose exhaustion is the manufactured exhaustion of labour performed under another's clock. Of that day's yield she keeps nothing; it was captured at the source, and no beauty she might have found in her toil would have returned it to her, which is why the paper refuses to tell her to find beauty in the toil. But the evening is not sold. The field she builds at home, with the one she loves, in the few unspacious hours that remain, yields a value of a different register, the real value that gathers in the relational field, and that yield is not captured, because it was never in the register the capturing reaches. The owner takes the day's holonomy; he cannot take the evening's, because the evening's accumulates in a coupling he has no purchase on, in a field whose gain returns to those who build it. This is why the cultivation matters most to her and not least: it is the one accumulation of her day that comes home with her, the one yield not siphoned, and to make her evening's repetition the gaining kind is to secure the only wealth the structure leaves in her hands. The structural injustice is real and the paper has named it; and within it, undefeated by it, the evening field remains hers to build, its yield the one the capture cannot reach.
\end{casebox}

\subsection{Aesthetic education and structural change}
\label{sub:pe-structural-change}

A question remains that the account of the structural remainder raised but did not settle: whether the aesthetic education of the everyday bears at all on the structural change that the structural remainder requires, or whether it is simply silent on it, a cultivation of the private field that leaves the public structure untouched. The honest answer is mixed, and stating it precisely guards against two errors. The first error would claim too much, that the cultivation of generative fields is itself a form of structural change, that to build accumulating fields in the everyday is to transform the conditions that capture their yield. This is false, and the paper has insisted on its falsity: the capture of the yield is a structural condition, and no quantity of beautiful field-building changes who owns the cycle's yield, which is altered only by a change in the structure of ownership. To claim that aesthetic education is structural change would be to aestheticize across the very line the paper has refused to cross. The second error would claim too little, that aesthetic education is wholly irrelevant to structural change, a private cultivation with no bearing on the public transformation. This too is false, and seeing why is the point of the subsection.

Aesthetic education bears on structural change in one specific and limited way, through the kind of subject it cultivates. A person practised in building generative fields, in producing meaning that is not borrowed, in perceiving who is rendered fuel within an accumulating field, is a person with capacities that structural change requires, the capacity to produce and not only consume meaning, the capacity to perceive injustice within apparent flourishing, the capacity to build the generative relations on which any reconstructed structure would depend. The cultivation does not itself change the structure, but it cultivates the subjects who could, and a structural change carried out by subjects incapable of building generative fields would rebuild the captured structure in a new form, as the series has noted in its analyses of how new forms of decentralised production are recaptured by the orders they were meant to escape. The bearing of aesthetic education on structural change is therefore neither identity nor irrelevance but cultivation of the agents: it does not change the structure, and it does form the kind of subject without whom a change in the structure would only reproduce the capture under a new name. This is a limited bearing, and the paper claims no more, but it is not nothing, and it places the aesthetic education of the everyday in a relation to the structural question that is honest about both its limits and its contribution, the limit that it cannot change the structure and the contribution that it forms the subjects any genuine change would require.

\subsection{Those who most need it}
\label{sub:pe-who-needs}

The section closes by reversing the charge entirely. The charge held that aesthetic education is for the leisured and irrelevant to the dispossessed; the account establishes that it is the dispossessed who most need it. Those whose days are full of owned, exhausting, captured repetition are exactly those for whom the difference between the wheel and the spiral is most consequential, because their repetition is most in danger of being the wheel, most pressed toward the mechanical and the deadening by the conditions that own it. For them, the capacity to make their repetition the gaining kind, to turn even the owned and exhausting day so that the field they build with those they love accumulates rather than flattens, is not a luxury but a defence, the one accumulation available to them when so much is captured, the real value that gathers in the relational field and cannot be siphoned because it was never in the register the siphoning reaches. The aesthetic education of the everyday is, for the dispossessed, the cultivation of the one cycle whose yield returns to them, the field they build between themselves and those they love, and that, far from being a luxury of the leisured, is the most necessary cultivation of all for those from whom every other yield is taken. The charge inverts: it is not that the dispossessed cannot afford aesthetic education, but that they, of all people, cannot afford to be without it.
